\documentclass[11pt]{article}
\usepackage[a4paper,margin=25mm]{geometry}
\usepackage{amsmath,amssymb,amsthm,mathtools,booktabs,hyperref,xurl}
\hypersetup{hidelinks}
\setlength{\emergencystretch}{2em}
\newtheorem{theorem}{Theorem}
\newtheorem{lemma}[theorem]{Lemma}
\newtheorem{proposition}[theorem]{Proposition}
\newtheorem{corollary}[theorem]{Corollary}
\newcommand{\Z}{\mathbb Z}\newcommand{\Q}{\mathbb Q}
\title{Both first-Pell signs have the same coding projection}
\author{\(\mathbb Z\mathbb Q\) Diophantine Project}
\date{23 November 2026}
\begin{document}
\maketitle

\begin{abstract}
The short-power source with the strengthened small-Pell margin no
longer needs its first-Pell sign square. The margin and the special
auxiliary parameter exclude the zero branch by a five-adic divisibility
obstruction. For the remaining branches, the old congruences force
the absolute power-of-two relation. An alternating binomial truncation
recovers central-binomial divisibility even when the original quotient
is negative. The auxiliary power index and its trace sign are then
recovered with either sign of the auxiliary parameter. Positive source
witnesses can consequently be reconstructed for the same coding prefix.
This deletes one private coordinate and one quadratic residual on the
two short-power routes. The equivalence is a statement about existence
and coding projection; it does not assert that every surviving auxiliary
coordinate extends to the old sign-selected zero set.
\end{abstract}

\section{Retained source and precise claim}

Use the complete source of \cite{SignedQ}, with its strengthened margin
and the cofinite instance range \(x>32z\). The product simplification
of \cite{Product} may be applied independently. Set
\[
 L=5^{60},\quad T_b=b^5-1,\quad
 A_0=\frac{b^9+8b^5+b-8}{2},\quad A=A_0+5T_bs,
 \quad m_*=2bA-b^2-1.
\]
Retain integer inputs \(b,w,s,c,h,\gamma,J,\delta,\alpha_*\),
the nonzero reciprocal integer input \(i\), and all prefix inputs.
The reciprocal for \(i\) is reconstructed only after \(c\ne0\),
as explained below. The auxiliary power relations are
\begin{equation}\label{power}
 C_*=L+\delta(A-1),\qquad
 D_*^2-(A^2-1)C_*^2=1,\qquad
 q=D_*-(A-b)C_*-\alpha_*m_*.
\end{equation}
The retained binomial relations have
\begin{equation}\label{bin-def}
 Y=sM,\quad B=2r+1,\quad
 a=(wM+1)rY,\quad p=2wr^2Y^2M,
 \quad k=r+1+h(p-1)
\end{equation}
and
\begin{align}
 \tau^2-(p^2-1)k^2&=1,\label{small}\\
 \eta-C_*^2&>0,&\eta&=k^2-4(c-kY)^2,\label{margin}\\
 d&=bw+(a-2)c+(4a-5)\gamma,\label{linear}\\
 d^2-(a^2-1)c^2&=1,\label{first}\\
 F^2-(a^2-1)i^2c^4&=1,\label{second}\\
 G&=a+F^2(c+1-a),& H&=c+F^2(B-c)+cFJ,\label{crt}\\
 I^2-(G^2-1)H^2&=1.\label{third}
\end{align}
All Pell trace witnesses are initially rational. Delete the sole
equation for the private rational coordinate \(\zeta\):
\begin{equation}\label{deleted}
                 \zeta^2=2\bigl(ad+(a^2-1)c+1\bigr).
\end{equation}

\begin{theorem}\label{main}
Every rational solution of this modified source satisfies
\[
 b\text{ is a positive power of two},\qquad
 M\mid\binom{2r}{r},\qquad q=b^L.
\]
It has an old sign-selected source witness for the same coding prefix
and instance, after the binomial and auxiliary-power witnesses are
chosen anew. Conversely every old source witness projects to a new
one by forgetting \(\zeta\). Thus deleting \eqref{deleted} preserves
the instance projection and rational undecidability transfer under
the same retained arithmetic-interface and coding premises.
\end{theorem}

Use the polynomial Pell sequences \(X_j(t),Y_j(t)\), with
\(X_0=1,X_1=t,Y_0=0,Y_1=1\) and recurrence
\(Z_{j+2}=2tZ_{j+1}-Z_j\). In particular
\[
 X_j(-t)=(-1)^jX_j(t),\qquad
 Y_j(-t)=(-1)^{j-1}Y_j(t).
\]
We use the fully signed index theorem of \cite{CRT}, which concludes
\(\lvert c\rvert=Y_B(\lvert a\rvert)\) and
\(\lvert d\rvert=X_B(\lvert a\rvert)\) without any sign square.

\section{Reconstruction before the first Pell index}

The preliminary power argument of \cite{Power} uses neither
\eqref{deleted} nor \(s\ne0\). For integer \(b\ge2\) and signed
integer \(s\), it proves
\[
 A_0,A,m_*\in\Z,\quad A\equiv1\pmod5,\quad |A|\ge2,
 \quad 0<A_0<b^{16},
 \quad m_*=T_b(b^5+8b+1+10bs).
\]
Indeed parity makes \(A_0\) integral, and reduction modulo five
gives \(A\equiv1\). The factorization excludes \(A=1\), since
it would imply \(T_b\mid-(b-1)^2\); the congruence excludes
\(A=0,-1\). Equations \eqref{power} then make \(C_*,D_*,q\)
integral. Pell classification permits index zero at this stage and
gives \(C_*=\pm Y_\nu(A)\). Its congruence modulo five makes
\(5\mid\nu\). The identities
\begin{equation}\label{power-rec}
 \begin{split}
 R_j&=X_j(A)-(A-b)Y_j(A),\qquad R_0=1,\\
 R_{j+1}&=bR_j+m_*Y_j(A),\\
 R_j\bigl(X_j(A)+(A-b)Y_j(A)\bigr)&=1+m_*Y_j(A)^2
 \end{split}
\end{equation}
make all four trace/coordinate signs congruent to
\(\pm b^{\pm\nu}\equiv\pm1\pmod{T_b}\). Hence
\[
 q\equiv\pm1\pmod{T_b},\qquad
 \lambda=(q^4-1)/T_b\in\Z.
\]
This argument does not invoke a sign of \(w,s,q,A\).

The signed-prefix lemmas in \cite{SignedQ} use only this integral
reconstruction, the unchanged positive prefix inputs, and the slack
inequality \(e\ell g^2<(b-xy)q^2\). With \(Q=|q|\), they give
\begin{equation}\label{prefix}
 b>xy\ge x>32z\ge32,\quad Q\ge b^5-2\ge b^4,
 \quad M=Q^{32},\quad r\ge M\ge b^{128}>2^{640}>L.
\end{equation}
In particular \(r,M\) are positive integers, \(b<Q\),
\(M\ge4b\), and \(M\ge16\). None of these conclusions uses
the first Pell sign square or any nonzero assertion about \(s\).

\begin{lemma}[The strengthened margin excludes zero]\label{nonzero}
The relations reconstructed so far imply \(s\ne0\).
\end{lemma}
\begin{proof}
Suppose \(s=0\). Then \(Y=a=p=0\), while \(k=r+1-h\)
is an integer. The strict margin implies \(k\ne0\), and
\eqref{small} says \(\tau^2+k^2=1\). Thus \(k=\pm1\),
\(\tau=0\). Since \(c,C_*\) are integers, the inequality
\(1-4c^2-C_*^2>0\) gives \(c=C_*=0\). Now \(A=A_0\),
so \eqref{power} forces \(A_0-1\mid L=5^{60}\).

For \(N_j(b)=1+b+\cdots+b^{j-1}\), factor
\begin{equation}\label{factor}
 A_0-1=(b-1)H_0(b),\qquad
 H_0(b)=\frac{N_9(b)+8N_5(b)+1}{2}.
\end{equation}
The numerator is even for every integer \(b\). Also \(H_0(1)=25\),
and its positive coefficients imply \(H_0(b)>25\) for \(b\ge2\).
If \(A_0-1\mid5^{60}\), then its positive factor \(b-1\)
equals \(5^a\) for some integer \(a\ge0\). Since \(b>32\),
one has \(a\ge3\). Thus \(b\equiv1\pmod{125}\), and
\eqref{factor}, with two invertible modulo 125, gives
\[
                        H_0(b)\equiv25\pmod{125}.
\]
Consequently \(v_5(H_0(b))=2\). A positive divisor of \(5^{60}\)
with that valuation is exactly 25, contradicting \(H_0(b)>25\).
\end{proof}

The margin now gives \(k\ne0\),
\begin{equation}\label{half}
 \left|\frac ck-Y\right|<\frac12,\qquad
 |c|>|k|(|Y|-1/2)>|k|>0,
 \qquad |C_*|<|k|<|c|,
\end{equation}
because \(|Y|\ge M\ge16\). This also justifies the retained
reciprocal reconstruction: the row \(vE=c^2\) forces \(v\ne0\),
and its integer predicate supplies \(i=1/v\in\Z\setminus\{0\}\).
The ordinary-\(abc\) chart additionally retains its explicit
\(vi=1\) row. No new reciprocal or positivity input is introduced.

\section{Both signs of the first trace}

\begin{lemma}\label{signed-binomial}
The retained binomial relations imply
\[
 |bw|=2^{2r+1},\qquad
 \operatorname{sgn}(c)=\operatorname{sgn}(w),\qquad
 \operatorname{sgn}(s)=\operatorname{sgn}(w)^{r+1}.
\]
If \(\sigma=\operatorname{sgn}(w)\), \(U=|w|M\),
\(Z=|s|M\), \(a_1=|a|\), \(P=|p|\), then
\begin{gather}
 C:=|c|=Y_{2r+1}(a_1),\qquad K:=|k|=Y_{r+1}(P),\label{values}\\
 a_1=rZ(U+\sigma),\qquad P=2r^2Z^2U,\qquad
 \left|\frac CK-Z\right|<\frac12,\label{magnitudes}\\
 \rho=\left(\frac{(U+\sigma)^2}{U}\right)^r,\qquad
 0<\rho-\frac CK<\frac{\rho}{Z^2(U-1)^2},\label{ratio}\\
 Z>\frac78(U-2)^r,\qquad a_1>2U4^r.\label{rough-growth}
\end{gather}
\end{lemma}
\begin{proof}
The preliminary arguments of \cite[Section 3]{Signed} use no sign
selection. For clarity, if \(w=0\), then \(p=0\), \(k=\pm1\),
and \eqref{half} forces \(c=kY\). The first norm would then give
\(d^2=(rY^2)^2-Y^2+1\), strictly between the consecutive squares
\((rY^2-1)^2\) and \((rY^2)^2\). Here \(d\in\Z\) by
\eqref{linear}. Thus \(w\ne0\), so \(U,Z\ge16\).

The expressions in \eqref{magnitudes} give \(a_1\ge2\),
\(P-1>5r+1\), \(2a_1<P^2\), and \(P>a_1\).
The congruence for \(k\), modulo \(|p-1|=P-\sigma\), gives
\(|k|\ge r+1\); hence \eqref{half} gives \(|c|>2r+1=B\).
The one-quotient index theorem of \cite{CRT}, with \eqref{first}
through \eqref{third}, yields
\[
                C=Y_B(a_1),\qquad |d|=X_B(a_1).
\]
The small norm makes \(\tau\) integral and writes
\(|k|=Y_\ell(P)\) for \(\ell\ge1\). The bounds
\(|k|<C<(2a_1)^{2r}<P^{4r}\) force \(\ell\le4r\).
Reducing \(k\equiv r+1\pmod{p-1}\), using
\(Y_\ell(P)\equiv\ell\pmod{P-1}\) or
\(Y_\ell(P)\equiv(-1)^{\ell-1}\ell\pmod{P+1}\),
now forces
\[
                         \ell=r+1,\qquad k=\sigma^rK.
\]
The modulus is greater than \(5r+1\), so no other signed residue
in the displayed index range can occur. Equation \eqref{half}
therefore gives \eqref{magnitudes} and
\(\operatorname{sgn}(c)=\operatorname{sgn}(s)\sigma^r\).

For completeness, put \(f_j(t)=Y_j(t)/(2t)^{j-1}\) for \(t\ge2\).
The normalized Pell lemma of \cite{Binomial} gives
\(1-(j-2)/(4t^2)\le f_j(t)\le1\), monotonicity in \(t\),
and strict decrease in \(j\) from index two. Thus
\[
 \frac CK=\rho\frac{f_{2r+1}(a_1)}{f_{r+1}(P)}<\rho,
 \qquad
 f_{2r+1}(a_1)>1-\frac1{2Z^2(U-1)^2}.
\]
These imply \eqref{ratio}. Since \(\rho>(U-2)^r\),
\eqref{magnitudes} gives
\[
 Z>\frac{15}{16}(U-2)^r-\frac12>\frac78(U-2)^r.
\]
Multiplying by \(r(U-1)\) proves \eqref{rough-growth}; explicitly,
\[
 \frac78r\frac{U-1}{U}\left(\frac{U-2}{4}\right)^r
 \ge\frac78\cdot\frac{15}{16}\cdot\frac{14}{4}
 =\frac{735}{256}>2.
\]

It remains to handle both independent first-Pell signs. Since \(B\)
is odd, write
\[
              d=\varepsilon X_B(a),\qquad
              c=\kappa Y_B(a),\qquad
              \varepsilon,\kappa\in\{1,-1\}.
\]
Here \(Y_B(a)=C>0\), irrespective of the sign of \(a\).
For \(m_0=4a-5\), the polynomial identities
\[
 \begin{split}
 R_j^-&=X_j(a)-(a-2)Y_j(a),\quad
 R_j^+=X_j(a)+(a-2)Y_j(a),\\
 R_{j+1}^-&=2R_j^-+m_0Y_j(a),\quad R_0^-=1,\qquad
 R_j^-R_j^+=1+m_0Y_j(a)^2
 \end{split}
\]
give \(R_B^-\equiv2^B\) and \(2^BR_B^+\equiv1\pmod{m_0}\).
The input bound \(b\le2M\) gives \(|bw|\le2U\), and
\eqref{rough-growth} gives
\begin{equation}\label{m0-bound}
 |m_0|\ge4a_1-5>8U4^r-5>4U4^r+1.
\end{equation}
If \(\kappa=-\varepsilon\), the linear row gives
\(bw2^B\equiv\varepsilon\pmod{m_0}\). This is impossible:
\[
 0<|bw2^B-\varepsilon|\le4U4^r+1<|m_0|.
\]
The first inequality follows from \(|bw|\ge2\).
Hence \(\kappa=\varepsilon\). The direct congruence and
\[
 |bw-\varepsilon2^B|\le2U+2\cdot4^r<|m_0|
\]
force \(bw=\varepsilon2^B\). Since \(b>0\), this means
\(\varepsilon=\kappa=\sigma\). The signs of \(c\) and \(s\)
follow from their already established relation.
\end{proof}

\section{Alternating truncation recovers divisibility}

\begin{lemma}\label{rounding}
With the notation of Lemma~\ref{signed-binomial},
\begin{equation}\label{truncation}
 Z=\sum_{j=0}^r\binom{2r}{j}\sigma^jU^{r-j}.
\end{equation}
In particular \(M\mid\binom{2r}{r}\).
\end{lemma}
\begin{proof}
Call the integer on the right of \eqref{truncation} \(T_\sigma\).
Expanding \(\rho=(U+\sigma)^{2r}/U^r\) gives
\[
 |\rho-T_\sigma|
 \le\sum_{j=r+1}^{2r}\binom{2r}{j}U^{r-j}
 <\frac{4^r/U}{1-1/U}.
\]
The exact absolute power relation gives
\(U=2\cdot4^rM/b\). Since \(M\ge4b\) and \(U\ge16\),
\begin{equation}\label{tail}
 |\rho-T_\sigma|<\frac{b}{2M}\frac{16}{15}
 \le\frac{2}{15}<\frac14.
\end{equation}
This estimate holds for both \(\sigma=1\) and \(\sigma=-1\).

Put \(E=\rho-C/K>0\) and \(D=Z^2(U-1)^2\). Equation
\eqref{ratio} says \(E<\rho/D=(C/K+E)/D\), so
\[
 E<\frac{C/K}{D-1}<\frac{Z+1/2}{Z^2(U-1)^2-1}<\frac14.
\]
The last inequality follows from \(U,Z\ge16\):
\(Z^2(U-1)^2>4Z+3\). Consequently
\[
 |Z-T_\sigma|
 \le|Z-C/K|+|C/K-\rho|+|\rho-T_\sigma|
 <\frac12+\frac14+\frac14=1.
\]
Both numbers are integers, so they are equal. As \(M\mid U\)
and \(M\mid Z\), reduction modulo \(M\) gives
\(M\mid\sigma^r\binom{2r}{r}\), which is the desired divisibility.
\end{proof}

This is the point where the strengthened prefix bound matters:
the conclusion uses \(M\ge4b\), not merely \(b\le2M\).
For negative \(w\), the quotient \(s\) need not be positive,
and the alternating formula must not be replaced by the old positive
truncation while keeping the same auxiliary coordinates.

\section{Power recovery with a signed auxiliary parameter}

Let \(S_1=|s|\), retaining the notation \(S,T\) of the coding
prefix separately. From \(|bw|=2^{2r+1}\) and \(b\le2M\),
\(U\ge4^r\). Bernoulli's inequality gives
\((1-2/U)^r\ge1-2r/U\ge1/2\). Lemma~\ref{signed-binomial}
therefore implies
\begin{equation}\label{growth}
 S_1>\frac7{16}\frac{U^r}{M}
 \ge\frac7{16}U^{r-1}>U^{r-2},\qquad
 |a|<2rS_1U^2<S_1U^3<S_1^2.
\end{equation}
Here \(M\le U\), \(2r<U\), and \(r\ge M>2^{640}\).
Moreover
\[
 S_1>U\ge4^r>7r+2>L+6r+2,\qquad
 S_1>U\ge M\ge b^{128}>A_0.
\]
For \(s>0\), the definition gives \(A>S_1\); for \(s<0\),
it gives \(A=A_0-5T_bS_1<-S_1\). Consequently
\begin{equation}\label{signed-A}
 |A|>S_1,\qquad |A-1|>L+6r,\qquad |m_*|>S_1.
\end{equation}
For the last assertion, if \(A<0\) it is immediate. If \(A>0\),
then \(m_*>2bS_1-b^2-1>S_1\), using \(S_1>b^{16}\).

Pell classification in \eqref{power} writes
\(C_*=\epsilon_*Y_\nu(A)\), \(D_*=\tau_*X_\nu(A)\),
with \(\nu\ge0\) and independent signs. If \(\nu=0\), the
congruence for \(C_*\) would make \(|A-1|\mid L\), contrary
to \eqref{signed-A}. For \(\nu\ge1\), elementary Pell growth,
\eqref{half}, and \eqref{growth} give
\[
 |A|^{\nu-1}\le|C_*|<|c|
 =Y_{2r+1}(|a|)<(2|a|)^{2r}<S_1^{6r}.
\]
Thus \(\nu\le6r\). Reduction modulo the positive modulus
\(|A-1|\) gives \(\epsilon_*\nu\equiv L\), because
\(A\equiv1\pmod{|A-1|}\) even when \(A<0\). Since
\(|A-1|>L+6r\), this forces \(\epsilon_*=1\) and \(\nu=L\).
In particular
\[
                  C_*=Y_L(A)>0,
\]
where the positivity also uses that \(L\) is odd.

Put \(K_*=b^L>1\). The polynomial identities \eqref{power-rec}
are valid with either sign of \(A\) and \(m_*\). They give
\[
 \begin{array}{ll}
 D_*=X_L(A):&q\equiv K_*\pmod{|m_*|},\\
 D_*=-X_L(A):&qK_*\equiv-1\pmod{|m_*|}.
 \end{array}
\]
The inverse is legitimate since \(m_*\equiv-1\pmod b\).
Since \(b<Q=M^{1/32}\le U^{1/32}\) and \(L<r\),
\[
 QK_*+1\le2QK_*\le2U^{(6r+1)/32}
 <U^{r-2}<S_1<|m_*|.
\]
For the strict middle inequality use \(U>2\), \(r\ge4\), and
\((6r+1)/32+1\le r-2\). In the negative trace case,
\(qK_*+1\) is a nonzero integer, since \(K_*>1\); its absolute
value is at most \(QK_*+1<|m_*|\). It cannot be a multiple
of \(|m_*|\). Thus \(D_*=X_L(A)\). Now
\(|q-K_*|\le Q+K_*\le QK_*+1<|m_*|\), so the direct
congruence forces \(q=K_*=b^L\).

\section{Rebuilding positive witnesses and the exact deletion}

\begin{proof}[Completion of Theorem~\ref{main}]
Lemma~\ref{signed-binomial} makes the positive integer \(b\) a
divisor of \(2^{2r+1}\), hence a positive power of two.
Lemma~\ref{rounding} supplies \(M\mid\binom{2r}{r}\), and
the preceding section supplies \(q=b^L\). Keep the coding-prefix
inputs, including \(b,q,r,M\), fixed. The positive binomial
certificate of \cite[Theorem 1]{Binomial} reconstructs positive
\(w,s,c,h,\gamma,i\) and the canonical first and small Pell
witnesses for exactly this \((r,M,b)\). The positive-completeness
proposition of \cite{CRT} supplies a positive one-quotient witness
\(J\). These choices need not agree with the original signed
auxiliary coordinates.

The auxiliary power completion in \cite{Power} supplies positive
\(\delta,\alpha_*\) and \(C_*=Y_L(A),D_*=X_L(A)\) for
the newly chosen positive \(s\). The absorbed-margin completeness
proof in \cite{SignedQ} applies with the same prefix bounds:
its canonical integer margin satisfies \(\eta\ge k>C_*^2\).
Thus the old positivity representatives can be chosen anew.
Finally the first-Pell half-trace identity gives
\(\zeta=2X_{r+1}(a)\), restoring \eqref{deleted}.

Every required integer input can therefore be chosen positive.
This permits the same thirteen-input general and ordinary-\(abc\)
interfaces to be reconstructed under their existing completeness
hypotheses, including the nonzero-input condition of the latter.
No particular old interface witness is required to survive this
reconstruction. Conversely, forgetting \(\zeta\) in any old
witness leaves every new row satisfied. This proves precisely the
claimed equivalence on coding-prefix and instance projections.
\end{proof}

The implemented half-trace residual is
\[
                    \zeta^2-2ad-2aA_c+2c-2,
                    \qquad A_c=ac.
\]
Its private coordinate \(\zeta\) occurs in no other residual.
Deleting this coordinate and row preserves the quadratic degree of
all remaining residuals; their real sum of squares still has degree
exactly four. This is one coordinate and one row on each of the two
short-power routes. The specified route still has its long fixed-power
circuit and does not have this strengthened auxiliary margin, so
no sign-square deletion is asserted there.

The complete predecessor is the general aggregate-spacing table in
\nolinkurl{undecidability/aggregate_spacing_v100/}, or the ordinary-\(abc\)
product table in \nolinkurl{undecidability/pell_product_v99/}.
The exact counts are
\begin{center}
\begin{tabular}{lrrrr}
\toprule
Interface & Old variables & New variables & Old rows & New rows\\
\midrule
General & 971 & 970 & 826 & 825\\
Ordinary \(abc\) & 696 & 695 & 591 & 590\\
\bottomrule
\end{tabular}
\end{center}
The independent full-table checker verifies the sole old occurrence of
\texttt{pell\_d\_positive\_square}, the exact deleted
\texttt{jones1980.pell\_positive\_root} row, and every retained row.
All thirteen integer-input interfaces remain. Historical metadata that
asserted positivity of the actual \(w,s\), or presence of a half-trace
square, is moved to predecessor history. The active metadata records
the signed conclusions and projection-only equivalence above.
The two tables and their certificates are reproduced with
\begin{verbatim}
python undecidability/pell_signed_branch_v101/build_signed_branch.py
python undecidability/pell_signed_branch_v101/verify_signed_branch.py
\end{verbatim}

The independent script \texttt{verify\_signed\_branch\_lemmas.py}
checks the parameter factorization, five-adic obstruction, signed
recurrence identities, alternating truncations, normalization and
modulus estimates, and exact signed binomial fixtures. It imports
no table builder. Its certificate records the scope of finite checks;
the quantified projection equivalence is the proof above.
In particular, it checks all 58 divisor candidates
\(b-1=5^a\), \(3\le a\le60\), after the proved zero-branch
reduction. Its ten large exact signed-binomial fixtures check the
first and small norms, integer congruence quotients, old strict margin,
and rounding estimates; they do not claim to instantiate the entire
coding prefix or the enormous third-Pell completion. The negative
\(w\) fixtures have negative right side in \eqref{deleted}, illustrating
why an identical-coordinate lift must not be claimed.

\begin{thebibliography}{9}
\small\setlength{\itemsep}{2pt}
\bibitem{Binomial} \(\mathbb Z\mathbb Q\) Diophantine Project,
\emph{An exact Pell certificate for a power of two and central-binomial
divisibility}.\par
\nolinkurl{undecidability/papers/binomial_pell_certificate.tex}.
\bibitem{Signed} \(\mathbb Z\mathbb Q\) Diophantine Project,
\emph{Signed binomial witnesses force their own positivity and save
five quadratic coordinates}.\par
\nolinkurl{undecidability/source_repair_v93/signed_binomial_positivity.tex}.
\bibitem{CRT} \(\mathbb Z\mathbb Q\) Diophantine Project,
\emph{A single integer quotient for the third Pell certificate}.\par
\nolinkurl{undecidability/pell_crt_v95/one_quotient_pell.tex}.
\bibitem{Power} \(\mathbb Z\mathbb Q\) Diophantine Project,
\emph{A short Pell replacement for the fixed Jones power}.\par
\nolinkurl{undecidability/pell_power_v97/pell_power_replacement.tex}.
\bibitem{SignedQ} \(\mathbb Z\mathbb Q\) Diophantine Project,
\emph{Signed power recovery and an absorbed Pell bound save six
quadratic coordinates}.\par
\nolinkurl{undecidability/pell_signed_q_v98/signed_q_and_absorbed_bound.tex}.
\bibitem{Product} \(\mathbb Z\mathbb Q\) Diophantine Project,
\emph{One product identity removes a quadratic coordinate}.\par
\nolinkurl{undecidability/pell_product_v99/pell_product_elimination.tex}.
\end{thebibliography}
\end{document}
